%!TEX root = ../QuantumNoiseAndDecoherence.tex
% ──────────────────────────────────────────────────────────────
%  Lecture 1 — Recapitulation of Quantum Mechanics
% ──────────────────────────────────────────────────────────────

\section{Recapitulation of quantum mechanics}\label{sec:recap-qm}

\subsection{Why quantum noise and decoherence?}

You have all met quantum mechanics before---it is a prerequisite for
this course.  But you have probably only encountered the quantum
mechanics of \emph{closed} systems: free particles, harmonic
oscillators, and other idealised setups that do not interact with
their environment.  The objective of this course is to go beyond this
and study the quantum mechanics of \emph{open} systems---quantum
systems composed of many subsystems.

In contrast to classical physics, quantum information \emph{cannot be
copied}: the \textbf{no-cloning theorem} prevents us from duplicating
quantum states.  If we want to gain information about a closed
quantum system, we have to interact with it, and this interaction
unavoidably has some back-action on the system itself.  \emph{The act
of measuring disturbs.}  This is the defining feature of quantum
mechanics.

\begin{intuition}[Measurement disturbance in everyday life]
  Measurement back-action is not unique to quantum mechanics.  Two
  classical examples:
  \begin{itemize}
    \item \textbf{Stock markets:} acting on information about a stock
      creates disturbance, since other traders observe and react to your
      trades.
    \item \textbf{Observer effect in social psychology:} subjects who
      know they are being observed behave differently.
  \end{itemize}
  In classical physics, these are consequences of approximations.  In
  quantum mechanics, measurement disturbance is \emph{fundamental}---it
  is built into the theory itself.
\end{intuition}

By the end of this course, we aim to understand this statement
quantitatively, and also to explain \emph{why quantum systems look
classical} when we observe them with our everyday measurement
apparatus.  The answer we propose: because of \textbf{quantum noise and
decoherence}, which are themselves a consequence of the fact that
measuring disturbs a quantum system.

\subsection{Course overview}

The course covers the following topics:
\begin{enumerate}[(1)]
  \item Recapitulation of quantum mechanics (complementarity and
    uncertainty, quantitative statements, entanglement, entropy); 1--2
    lectures.
  \item Complete positivity and dilation theory (quantum channels, POVM
    observables, state after general measurements, non-demolition
    measurements, Naimark, Stinespring); 1~lecture.
  \item Quantum many-particle systems (Fock space, quasi-free states,
    mean-field limits); 2~lectures.
  \item Continuous variables (covariance matrices, Gaussian states,
    Wigner distributions, examples from quantum optics, standard
    quantum limit, continuous variables from collective degrees of
    freedom); 2~lectures.
  \item Lindblad equations (physical background: quantum dynamics with
    dissipation and decoherence, Markov approximation, derivations:
    weak and strong coupling regime); 2~lectures.
  \item Decoherence sources in experiments (the basic mechanism: losing
    parts of the system, decoherence of Schr\"odinger cat states,
    classical limit); 1~lecture.
  \item Classical noise and stochastic calculus (basics of stochastic
    processes, counting processes, photodetection, factorial moments,
    Markov chains and processes, martingales, Langevin and
    Fokker--Planck equations); 2~lectures.
  \item Quantum stochastic calculus (discrete case: matrix product
    states, unravelling of Lindblad equations, quantum statistics of
    the output events, continuous quantum measurement, cavity
    input--output formalism).
  \item Feedback and control (classical and quantum feedback, Lie
    algebra techniques, controllability, linear quantum control
    systems); 1~lecture.
\end{enumerate}

\medskip
\noindent\textbf{Books/Notes:}
\begin{enumerate}[(1)]
  \item Reinhard~--- Introduction to QM + QI~\cite{Reinhard2012}
  \item B.~Davies~\cite{Davies76}
  \item M.~Wolf's notes~\cite{Wolf12}
  \item Gardiner and Zoller~\cite{GardinerZoller04}
  \item Gardiner~\cite{Gardiner09}
  \item Wiseman \& Milburn~\cite{WisemanMilburn10}
  \item Walls \& Milburn~\cite{WallsMilburn08}
\end{enumerate}

% ──────────────────────────────────────────────────────────────
\subsection{Kinematics of quantum mechanics}\label{sec:kinematics}

We divide physical experiments into two parts:

\begin{center}
\begin{tikzpicture}[
  block/.style={ellipse, draw=cgblue, fill=cgblue!10,
    minimum height=2em, minimum width=3.5cm,
    font=\sf\small, text=spacecadet, inner sep=5pt},
  arr/.style={-{Stealth[length=5pt]}, thick, cgblue}
]
  \node[block] (prep) {Preparation};
  \node[block, right=2cm of prep] (meas) {Measurement};
  \draw[arr] (prep) -- (meas);
\end{tikzpicture}
\end{center}

In classical physics there is no real need to talk of measurements.  In
statistical theories, such as quantum mechanics, all we can do is predict
probabilities of outcomes in statistical experiments.

\begin{definition}[Preparation]\label{def:preparation}
  A \emph{preparation} of a quantum system is the set of actions which
  determines all probability distributions of any possible measurement.
  Different preparations can lead to identical probability
  distributions~$\Rightarrow$ we introduce the concept of \emph{state},
  which specifies only the effect of preparation regardless of how it was
  performed.
\end{definition}

\begin{keyresult}[State]
  A \emph{state} describes all details relevant for subsequent
  measurement.
\end{keyresult}

Mathematically, the division of experiments into preparations and
measurements of observables is represented, in quantum mechanics, as
follows.

\subsubsection{Observables}
\emph{Observables} are hermitian elements of an algebra~$\alg$ (think:
matrices of side~$n$).

An \emph{algebra} is a set closed under multiplication and addition as
well as multiplication by scalars.  Each element $A \in \alg$ has an
adjoint $A^\adj \in \alg$.  An element $E \in \alg$ is \emph{positive}
if $E = A^\adj A$.

The algebra~$\alg$ is usually represented in terms of bounded linear
operators $\BH$ acting on a Hilbert space~$\hilb$.

\begin{example}[Qubit observables]\label{ex:qubit-obs}
  The set of observables for a \emph{qubit} are hermitian $2\times 2$
  matrices in
  \[
    \alg_n \cong \alg \cong \Mn{2}
    = \left\{
      \begin{pmatrix} a & b \\ c & d \end{pmatrix}
      \;\middle|\; a,b,c,d \in \C
    \right\}\,,
    \qquad
    \alg_h = \{ X \in \alg \mid X^\adj = X \}\,.
  \]
\end{example}

\subsubsection{Effects and measurements}\label{sec:effects}
Measurements themselves (not observables) are described by assigning to
each \emph{outcome} of a device an \emph{effect} $E \in \alg$: an
element satisfying
\begin{equation}\label{eq:effect}
  0 \leq E \leq \One\,.
\end{equation}
Here $A \leq B$ means $B - A = X^\adj X$ for some~$X$.

\begin{example}[Photon counter]\label{ex:photon-counter}
  A photon counter has two possible outcomes: no click or click.  The
  corresponding effects are:
  \begin{align}
    E_{\text{no click}} &= \ketbra{0}{0}\,,
    & E_{\text{click}} &= \One - \ketbra{0}{0}\,,
  \end{align}
  where $\ket{0}$ is the vacuum state.  Homodyne detection, by
  contrast, assigns to each outcome a projection onto a coherent
  state---an example of a POVM that does \emph{not} arise from a
  standard textbook observable.
\end{example}

\subsubsection{States}\label{sec:states}
A \emph{state}~$\omega$ on~$\alg$ is a positive normalised linear
functional on~$\alg$.  That is:
\begin{equation}\label{eq:state-def}
  \omega\colon \alg \to \C \text{ is linear},\quad
  \omega(X^\adj X) \geq 0,\quad
  \omega(\One) = 1\,.
\end{equation}

We think of states as assigning probabilities to outcomes of
measurements.  Each of the three axioms has a physical origin:
\begin{itemize}
  \item \textbf{Normalisation} $\omega(\One) = 1$: the identity is the
    effect corresponding to ``doing nothing whatsoever.''  Its
    probability must be~$1$ (probability is conserved).
  \item \textbf{Positivity} $\omega(X^\adj X) \geq 0$: if we put in a
    valid measurement (a positive operator), we must get out a valid
    probability (a non-negative number).
  \item \textbf{Linearity}: if we randomise between two preparations
    (e.g.\ flip a fair coin and perform preparation~$A$ on heads,
    preparation~$B$ on tails), the resulting state is a convex
    combination $\tfrac{1}{2}\omega_A + \tfrac{1}{2}\omega_B$.
    Linearity ensures this is consistent.
\end{itemize}

\begin{remark}
  States described as rays in Hilbert space are a \emph{subset} of all
  possible states: they are the pure states.  In this course we deal
  with the real world, where noise and imperfections force quantum
  systems into \emph{mixed} states, and we therefore need the full
  generality of the algebraic framework.
\end{remark}

\subsubsection{Probabilities and the Born rule}
The probability of outcome~$E$ measured on a state~$\omega$ is
\begin{equation}\label{eq:born-rule}
  \eqbox{p(E) = \omega(E)}\,.
\end{equation}

Label outcomes corresponding to observable~$A$ by index~$\alpha$.
We assign an effect $0 \leq E_\alpha \leq \One$ to each~$\alpha$ so
that $p_\alpha = \text{prob.\ of outcome } \alpha$ can be calculated
via
\begin{equation}\label{eq:prob-outcome}
  p_\alpha = \omega(E_\alpha)\,.
\end{equation}
In order that $\sum_\alpha p_\alpha = 1$, we require
\begin{equation}\label{eq:povm-resolution}
  \eqbox{\sum_\alpha E_\alpha = \One}\,.
\end{equation}

\begin{definition}[POVM]\label{def:povm}
  A collection $\{E_\alpha\}$ of effects such that $\sum_\alpha E_\alpha
  = \One$ is called a \emph{POVM} (positive operator-valued measure).
\end{definition}

\subsubsection{Connection to textbook observables}
The connection between textbook observables~$A$ and POVMs is via the
spectral decomposition:
\begin{equation}\label{eq:spectral}
  A = \sum_\alpha a_\alpha\, E_\alpha\,,
\end{equation}
where $a_\alpha$ are the eigenvalues and $E_\alpha$ are the spectral
projectors.  Note $\sum_\alpha E_\alpha = \One$.

POVMs are \emph{more general}: the $E_\alpha$ need not be orthogonal
projectors.  Every observable gives rise to a POVM (via its spectral
decomposition), but not every POVM comes from an observable.  This extra
generality is very useful: measurements such as homodyne and heterodyne
detection in quantum optics do \emph{not} correspond to measuring a
standard hermitian observable---they are inherently POVM measurements.

\begin{intuition}[Why POVMs matter]
  It is somewhat puzzling that textbook quantum mechanics emphasises
  ``observables correspond to hermitian operators,'' when in practice
  many real experiments (heterodyne detection, weak measurements, etc.)
  perform measurements that are strictly more general than projective
  measurements of a hermitian operator.  The POVM framework captures
  \emph{all} possible measurements in a unified way.
\end{intuition}

% ──────────────────────────────────────────────────────────────
\subsection{Density operators}\label{sec:density-ops}

For finite-dimensional systems, e.g.\ qubits, we can \emph{always}
represent states $\omega\colon \alg \to \C$ as \emph{density operators}
via
\begin{equation}\label{eq:density-op}
  \eqbox{\omega(A) = \tr(\rho\, A)}\,,
\end{equation}
where
\begin{equation}\label{eq:density-props}
  \tr(\rho) = 1 \qquad\text{and}\qquad \rho \geq 0\,.
\end{equation}
Linearity of the state is automatic, since the trace is linear.

\begin{remark}
  The identification of states with density operators only works
  reliably for finite-dimensional systems.  In infinite dimensions
  (e.g.\ the electromagnetic field), there exist operators whose trace
  diverges, and the set of states is strictly larger than the set of
  density operators.  This is why we emphasise the algebraic definition
  of states as functionals on observables---it remains valid in all
  dimensions.
\end{remark}

\begin{definition}[Pure state]\label{def:pure-state}
  A state~$\omega$ is \emph{pure} if it cannot be written as a
  nontrivial convex combination (i.e.\ with $p \notin \{0,1\}$):
  \begin{equation}\label{eq:pure-state-def}
    \omega = (1 - p)\,\omega' + p\,\omega''\,.
  \end{equation}
  For finite-dimensional systems, $\omega$ is pure if and only if
  \begin{equation}\label{eq:pure-density}
    \omega(\cdot) = \tr(\rho\,\cdot) \quad\text{with}\quad
    \rho = \ketbra{\psi}{\psi}\,,\quad
    \ket{\psi} \in \hilb\,.
  \end{equation}
  That is, a pure state corresponds to a rank-one projector---a
  ``ray'' in Hilbert space.
\end{definition}

\begin{example}[Qubit states and the Bloch sphere]\label{ex:bloch}
  All qubit states~$\omega$ may be represented as $2 \times 2$
  matrices $\rho \geq 0$, $\tr(\rho) = 1$.  The most general such
  matrix that satisfies $\tr(\rho)=1$ can be parameterised as:
  \begin{equation}\label{eq:bloch-rho}
    \rho = \frac{1}{2}
    \begin{pmatrix}
      1 + x_3 & x_1 - ix_2 \\
      x_1 + ix_2 & 1 - x_3
    \end{pmatrix}
    = \frac{\One + \vb{x}\cdot\paulivec}{2}\,,
  \end{equation}
  where $\paulivec = (\sx,\,\sy,\,\sz)$ are the Pauli matrices:
  \begin{equation}\label{eq:pauli}
    \sx = \begin{pmatrix} 0 & 1 \\ 1 & 0 \end{pmatrix},\quad
    \sy = \begin{pmatrix} 0 & -i \\ i & 0 \end{pmatrix},\quad
    \sz = \begin{pmatrix} 1 & 0 \\ 0 & -1 \end{pmatrix},\quad
    \vb{x} \in \R^3\,.
  \end{equation}
\end{example}

\begin{exercise}\label{ex:bloch-ball}
  Show that $\rho \geq 0 \;\Rightarrow\; \lvert\vb{x}\rvert^2 \leq 1$.

  \emph{Hint:} Compute the eigenvalues of~$\rho$ (they are $p$ and
  $1-p$ for some $p \in [0,1]$) and impose $\rho \geq 0$.
  Alternatively, use the constraint $\rho^2 = \rho$ for pure states to
  show that pure states satisfy $\lvert\vb{x}\rvert = 1$, and then
  argue that mixed states lie in the interior by convexity.
\end{exercise}

Thus $\vb{x}$ lies in the unit ball of radius~$1$ in~$\R^3$---the
\emph{Bloch ball}---and the density operator~$\rho$ is in one-to-one
correspondence with~$\vb{x}$.  Pure states live on the
\emph{surface} of the sphere ($\lvert\vb{x}\rvert = 1$); mixed states
occupy the interior.

\begin{center}
\begin{tikzpicture}[scale=1.3]
  % Bloch sphere
  \shade[ball color=cgblue!15, opacity=0.4] (0,0) circle (1.2);
  \draw[thick, cgblue!60] (0,0) circle (1.2);
  % Equator (ellipse for 3D effect)
  \draw[cgblue!40, dashed] (0,0) ellipse (1.2 and 0.4);
  % Axes
  \draw[qarrow] (0,0) -- (1.6,0) node[right, font=\small]{$x_1$};
  \draw[qarrow] (0,0) -- (0.6,0.35) node[right, font=\small]{$x_2$};
  \draw[qarrow] (0,0) -- (0,1.6) node[above, font=\small]{$x_3$};
  % South pole
  \draw[qarrow, spacecadet] (0,0) -- (0,-1.5);
  % Bloch vector
  \draw[thick, -{Stealth[length=5pt]}, munsell] (0,0) -- (0.5,0.9)
    node[above right, font=\small, text=munsell]{$\vb{x}$};
  \fill[munsell] (0.5,0.9) circle (1.5pt);
  % Labels
  \node[font=\footnotesize, text=spacecadet] at (1.2,-1.5)
    {pure: $|\vb{x}|=1$};
  \node[font=\footnotesize, text=spacecadet] at (-1.5,0)
    {mixed: $|\vb{x}|<1$};
\end{tikzpicture}
\end{center}

\begin{intuition}[The Bloch sphere]
  It is a remarkable fact that qubit states---abstract $2\times 2$
  positive matrices with unit trace---admit a geometric representation
  as points in a three-dimensional ball.  Operations on qubits
  (rotations, measurements, decoherence) all have clean geometrical
  interpretations in terms of this object.  Sadly, for qutrits ($3$-level
  systems) and beyond, no comparably simple geometric picture exists:
  the state space becomes a complicated convex body in $d^2-1$
  dimensions with ``horrible corners and curvy bits.''
\end{intuition}

\begin{example}[Effect measurement on a qubit]\label{ex:qubit-effect}
  Consider the effect
  \[
    E = \begin{pmatrix} 1 & 0 \\ 0 & 0 \end{pmatrix}
    = \tfrac{1}{2}(\One + \sz)\,.
  \]
  This corresponds to measuring spin-up in the $z$-direction.  The
  probability of this outcome is
  \[
    p(E) = \omega(E) = \tr(\rho\, E) = \tfrac{1}{2}(1 + x_3)\,.
  \]
  Geometrically, this probability is determined by the projection of the
  Bloch vector~$\vb{x}$ onto the $x_3$-axis.
\end{example}

% ──────────────────────────────────────────────────────────────
\subsection{Composite systems and entanglement}\label{sec:composite}

The observables for a composite quantum system~$AB$ with observable
algebras $\alg$ and~$\algB$, respectively, is
\begin{equation}\label{eq:composite-alg}
  \alg \tp \algB\,.
\end{equation}
The underlying Hilbert space for~$AB$ is a tensor product as well:
\begin{equation}\label{eq:composite-hilb}
  \hilb = \hilbA \tp \hilbB\,.
\end{equation}

The simplest states~$\omega$ of a composite system~$AB$ are
\emph{products}, i.e.\
\begin{equation}\label{eq:product-state}
  \omega(X \tp Y) = \omega_A(X)\cdot\omega_B(Y)\,,\qquad
  X \in \alg,\; Y \in \algB\,.
\end{equation}
Pure product states~$\omega$ have density operator
\begin{equation}\label{eq:product-density}
  \rho = \ketbra{\psi}{\psi} \tp \ketbra{\phi}{\phi}\,,\qquad
  \ket{\psi} \in \hilbA,\;\ket{\phi} \in \hilbB\,.
\end{equation}

States which are \emph{not} product are \emph{correlated}.

\subsubsection{Reduced states and partial trace}
The state of~$AB$, when restricted to observables acting nontrivially
only on~$A$, i.e.\ observables of the form $X_A \tp \One_B$, is called
the \emph{reduced state}~$\omega_A$.  The corresponding \emph{reduced
density operator} $\rho_A$ on~$\alg$ is defined by
\begin{equation}\label{eq:reduced-state}
  \omega_A(X) = \tr[\rho_A\, X] = \tr[\rho\,(X \tp \One)]
  \qquad \forall\, X \in \alg\,.
\end{equation}
The map
\begin{equation}\label{eq:partial-trace}
  \eqbox{\tr_B\colon \rho \mapsto \rho_A
  \coloneqq \sum_{j} (\One \tp \bra{j})\,\rho\,(\One \tp \ket{j})}\,,
\end{equation}
where $\{\ket{j}\}$ is an ONB of~$\hilbB$, is called the
\emph{partial trace}; $\rho_A$ is the \emph{marginal} of~$\rho$.

\begin{center}
\begin{tikzpicture}[
  sysbox/.style={rectangle, rounded corners=3pt, draw=cgblue,
    fill=cgblue!8, minimum height=2em, minimum width=1.8cm,
    font=\sf\small, text=spacecadet, inner sep=5pt},
  arr/.style={-{Stealth[length=5pt]}, thick, cgblue}
]
  % Composite system
  \node[sysbox, minimum width=4.5cm] (AB) {$\omega_{AB}$};
  \node[font=\footnotesize, text=spacecadet] at ([yshift=8pt]AB.north) {$AB$};
  % Subsystems inside
  \draw[dashed, cgblue!40] (AB.center) -- (AB.south);
  \draw[dashed, cgblue!40] (AB.center) -- (AB.north);
  % Arrow to reduced
  \node[sysbox, below right=1.2cm and 0.5cm of AB] (A) {$\omega_A$};
  \node[font=\footnotesize, text=spacecadet] at ([yshift=8pt]A.north) {$A$};
  \draw[arr] (AB.south east) -- (A.north west)
    node[midway, right, font=\footnotesize]{$\tr_B$};
  % Discard label
  \node[font=\footnotesize, text=spacecadet!60, below left=1.2cm and 0.2cm of AB]
    {discard $B$};
\end{tikzpicture}
\end{center}

% ──────────────────────────────────────────────────────────────
\subsection{Schmidt decomposition}\label{sec:schmidt}

\begin{proposition}[Schmidt decomposition]\label{prop:schmidt}
  For every vector $\ket{\psi} \in \hilbA \tp \hilbB$, with
  $d = \min\{\dim(\hilbA),\, \dim(\hilbB)\}$, there exists an ONB
  $\{\ket{g_j}\} \subset \hilbA$ and $\{\ket{f_j}\} \subset \hilbB$
  such that
  \begin{equation}\label{eq:schmidt}
    \eqbox{\ket{\psi}
      = \sum_{j=1}^{d} \sqrt{p_j}\;\ket{g_j} \tp \ket{f_j}}\,,\qquad
    p_j \geq 0,\quad
    \sum_j p_j = \norm{\ket{\psi}}^2\,.
  \end{equation}
  The $\sqrt{p_j}$ are called the \emph{Schmidt coefficients}; the
  number of nonzero~$p_j$ is the \emph{Schmidt rank}.
\end{proposition}

\begin{proof}
  Write $\ket{\psi}$ in a product basis:
  \[
    U \tp V \bigl(\ket{\psi}\bigr)
    = U \tp V \Bigl(\sum_{j,k} [\Psi]_{jk}\,\ket{j}\ket{k}\Bigr)
    = \sum_{j,j'}\sum_{k,k'} [U]_{jj'}\,[\Psi]_{j'k'}\,[V]_{k'k}\,
      \ket{j}\ket{k}\,.
  \]
  By the singular value decomposition (SVD), choose unitary $U,V$ such
  that $U\Psi V^\top = D$ is diagonal with non-negative entries.  Then
  \[
    \ket{\psi}
    = \sum_{j} [D]_{jj}\,\ket{g_j}\tp\ket{f_j}\,,
  \]
  where $\ket{g_j} = U^\adj\ket{j}$ and $\ket{f_j} = V^\adj\ket{j}$
  form orthonormal bases of~$\hilbA$ and~$\hilbB$ respectively.
  Setting $\sqrt{p_j} = [D]_{jj}$ gives the Schmidt form.
\end{proof}
